\documentclass[11pt]{article}
% ============================================================================
% preamble.tex -- shared style for every document in docs/
%
% Each document starts with:
%     \documentclass[11pt]{article}
%     % ============================================================================
% preamble.tex -- shared style for every document in docs/
%
% Each document starts with:
%     \documentclass[11pt]{article}
%     % ============================================================================
% preamble.tex -- shared style for every document in docs/
%
% Each document starts with:
%     \documentclass[11pt]{article}
%     \input{preamble}
%     \newcommand{\doctitle}{...}
%
% Team names live in ONE place (below). Edit them here and every document
% updates on the next build.
% ============================================================================

% ---------------------------------------------------------------- team
\newcommand{\projname}{FPGA int8 MLP Accelerator}
\newcommand{\repourl}{https://github.com/itsKannapi/fpga-mlp-accelerator}
\newcommand{\memberA}{Joey (\texttt{itsKannapi})}
\newcommand{\memberB}{Teammate~2 (TBD)}
\newcommand{\memberC}{Teammate~3 (TBD)}

% ---------------------------------------------------------------- packages
\usepackage[margin=1in]{geometry}
\usepackage[T1]{fontenc}
\usepackage{lmodern}
\usepackage{microtype}
\usepackage{amsmath}
\usepackage{booktabs}
\usepackage{tabularx}
\usepackage{longtable}
\usepackage{array}
\usepackage{enumitem}
\usepackage{xcolor}
\usepackage{listings}
\usepackage{tikz}
\usetikzlibrary{arrows.meta, positioning, fit, calc}
\usepackage[most]{tcolorbox}
\usepackage{fancyhdr}
\usepackage{titlesec}
\usepackage{needspace}
\usepackage[hidelinks]{hyperref}
\hypersetup{colorlinks=true, linkcolor=accent, urlcolor=accent, citecolor=accent}

% ---------------------------------------------------------------- colors
\definecolor{accent}{HTML}{1F4E79}
\definecolor{muted}{HTML}{5F6B7A}
\definecolor{rule}{HTML}{C9D1DA}
\definecolor{okc}{HTML}{2E7D32}
\definecolor{progc}{HTML}{B26A00}
\definecolor{todoc}{HTML}{8A8F98}
\definecolor{codebg}{HTML}{F5F7FA}

% ---------------------------------------------------------------- layout
\setcounter{tocdepth}{1}   % sections only, keeps the contents to a few lines
\setlength{\parindent}{0pt}
\setlength{\parskip}{0.55em}
\setlist{itemsep=0.15em, topsep=0.3em}
\renewcommand{\arraystretch}{1.18}

\titleformat{\section}{\Large\bfseries\color{accent}}{\thesection}{0.7em}{}
\titleformat{\subsection}{\large\bfseries}{\thesubsection}{0.6em}{}
\titleformat{\subsubsection}{\normalsize\bfseries}{\thesubsubsection}{0.5em}{}
\titlespacing*{\section}{0pt}{1.4em}{0.5em}

\pagestyle{fancy}
\fancyhf{}
\fancyhead[L]{\small\color{muted}\projname}
\fancyhead[R]{\small\color{muted}\doctitle}
\fancyfoot[C]{\small\color{muted}\thepage}
\renewcommand{\headrulewidth}{0.4pt}
\renewcommand{\headrule}{\color{rule}\hrule width\headwidth height\headrulewidth}

% ---------------------------------------------------------------- title block
\newcommand{\makedoctitle}[2]{%  #1 = subtitle, #2 = last-updated date
  \thispagestyle{plain}
  {\color{muted}\small\MakeUppercase{\projname}}\par\vspace{0.2em}
  {\Huge\bfseries\color{accent}\doctitle}\par\vspace{0.4em}
  {\large #1}\par\vspace{0.6em}
  {\small\color{muted}Team: \memberA, \memberB, \memberC\quad$\cdot$\quad
   Last updated #2\quad$\cdot$\quad\href{\repourl}{\texttt{\repourl}}}\par
  \vspace{0.4em}{\color{rule}\rule{\linewidth}{0.6pt}}\par\vspace{0.6em}
}

% ---------------------------------------------------------------- status tags
\newcommand{\tagbox}[2]{\tcbox[on line, colback=#1!12, colframe=#1!12,
  coltext=#1!70!black, boxsep=0pt, left=3pt, right=3pt, top=1.5pt, bottom=1.5pt,
  arc=2pt]{\footnotesize\bfseries #2}}
\newcommand{\done}{\tagbox{okc}{DONE}}
\newcommand{\inprog}{\tagbox{progc}{IN PROGRESS}}
\newcommand{\todo}{\tagbox{todoc}{TO DO}}
\newcommand{\stretchgoal}{\tagbox{accent}{STRETCH}}

% ---------------------------------------------------------------- callouts
\newtcolorbox{note}[1][Note]{colback=accent!5, colframe=accent!60,
  boxrule=0.6pt, arc=2pt, left=6pt, right=6pt, top=4pt, bottom=4pt,
  fonttitle=\bfseries\small, title=#1, coltitle=white}
\newtcolorbox{warn}[1][Watch out]{colback=progc!6, colframe=progc!70,
  boxrule=0.6pt, arc=2pt, left=6pt, right=6pt, top=4pt, bottom=4pt,
  fonttitle=\bfseries\small, title=#1, coltitle=white}

% ---------------------------------------------------------------- code
\lstdefinelanguage{SV}{
  morekeywords={module,endmodule,input,output,wire,reg,signed,always,posedge,
    negedge,if,else,begin,end,assign,parameter,localparam,generate,endgenerate,
    for,genvar,case,endcase,function,endfunction,integer,initial},
  sensitive=true, morecomment=[l]{//}, morecomment=[s]{/*}{*/},
  morestring=[b]", morestring=[b]`}
\lstset{
  basicstyle=\ttfamily\small, backgroundcolor=\color{codebg},
  keywordstyle=\color{accent}\bfseries, commentstyle=\color{muted}\itshape,
  frame=single, rulecolor=\color{rule}, framesep=4pt, xleftmargin=4pt,
  xrightmargin=4pt, columns=fullflexible, keepspaces=true, breaklines=true,
  showstringspaces=false, upquote=true, aboveskip=0.6em, belowskip=0.6em}
\newcommand{\cmd}[1]{\texttt{#1}}
\newcommand{\file}[1]{\texttt{#1}}

%     \newcommand{\doctitle}{...}
%
% Team names live in ONE place (below). Edit them here and every document
% updates on the next build.
% ============================================================================

% ---------------------------------------------------------------- team
\newcommand{\projname}{FPGA int8 MLP Accelerator}
\newcommand{\repourl}{https://github.com/itsKannapi/fpga-mlp-accelerator}
\newcommand{\memberA}{Joey (\texttt{itsKannapi})}
\newcommand{\memberB}{Teammate~2 (TBD)}
\newcommand{\memberC}{Teammate~3 (TBD)}

% ---------------------------------------------------------------- packages
\usepackage[margin=1in]{geometry}
\usepackage[T1]{fontenc}
\usepackage{lmodern}
\usepackage{microtype}
\usepackage{amsmath}
\usepackage{booktabs}
\usepackage{tabularx}
\usepackage{longtable}
\usepackage{array}
\usepackage{enumitem}
\usepackage{xcolor}
\usepackage{listings}
\usepackage{tikz}
\usetikzlibrary{arrows.meta, positioning, fit, calc}
\usepackage[most]{tcolorbox}
\usepackage{fancyhdr}
\usepackage{titlesec}
\usepackage{needspace}
\usepackage[hidelinks]{hyperref}
\hypersetup{colorlinks=true, linkcolor=accent, urlcolor=accent, citecolor=accent}

% ---------------------------------------------------------------- colors
\definecolor{accent}{HTML}{1F4E79}
\definecolor{muted}{HTML}{5F6B7A}
\definecolor{rule}{HTML}{C9D1DA}
\definecolor{okc}{HTML}{2E7D32}
\definecolor{progc}{HTML}{B26A00}
\definecolor{todoc}{HTML}{8A8F98}
\definecolor{codebg}{HTML}{F5F7FA}

% ---------------------------------------------------------------- layout
\setcounter{tocdepth}{1}   % sections only, keeps the contents to a few lines
\setlength{\parindent}{0pt}
\setlength{\parskip}{0.55em}
\setlist{itemsep=0.15em, topsep=0.3em}
\renewcommand{\arraystretch}{1.18}

\titleformat{\section}{\Large\bfseries\color{accent}}{\thesection}{0.7em}{}
\titleformat{\subsection}{\large\bfseries}{\thesubsection}{0.6em}{}
\titleformat{\subsubsection}{\normalsize\bfseries}{\thesubsubsection}{0.5em}{}
\titlespacing*{\section}{0pt}{1.4em}{0.5em}

\pagestyle{fancy}
\fancyhf{}
\fancyhead[L]{\small\color{muted}\projname}
\fancyhead[R]{\small\color{muted}\doctitle}
\fancyfoot[C]{\small\color{muted}\thepage}
\renewcommand{\headrulewidth}{0.4pt}
\renewcommand{\headrule}{\color{rule}\hrule width\headwidth height\headrulewidth}

% ---------------------------------------------------------------- title block
\newcommand{\makedoctitle}[2]{%  #1 = subtitle, #2 = last-updated date
  \thispagestyle{plain}
  {\color{muted}\small\MakeUppercase{\projname}}\par\vspace{0.2em}
  {\Huge\bfseries\color{accent}\doctitle}\par\vspace{0.4em}
  {\large #1}\par\vspace{0.6em}
  {\small\color{muted}Team: \memberA, \memberB, \memberC\quad$\cdot$\quad
   Last updated #2\quad$\cdot$\quad\href{\repourl}{\texttt{\repourl}}}\par
  \vspace{0.4em}{\color{rule}\rule{\linewidth}{0.6pt}}\par\vspace{0.6em}
}

% ---------------------------------------------------------------- status tags
\newcommand{\tagbox}[2]{\tcbox[on line, colback=#1!12, colframe=#1!12,
  coltext=#1!70!black, boxsep=0pt, left=3pt, right=3pt, top=1.5pt, bottom=1.5pt,
  arc=2pt]{\footnotesize\bfseries #2}}
\newcommand{\done}{\tagbox{okc}{DONE}}
\newcommand{\inprog}{\tagbox{progc}{IN PROGRESS}}
\newcommand{\todo}{\tagbox{todoc}{TO DO}}
\newcommand{\stretchgoal}{\tagbox{accent}{STRETCH}}

% ---------------------------------------------------------------- callouts
\newtcolorbox{note}[1][Note]{colback=accent!5, colframe=accent!60,
  boxrule=0.6pt, arc=2pt, left=6pt, right=6pt, top=4pt, bottom=4pt,
  fonttitle=\bfseries\small, title=#1, coltitle=white}
\newtcolorbox{warn}[1][Watch out]{colback=progc!6, colframe=progc!70,
  boxrule=0.6pt, arc=2pt, left=6pt, right=6pt, top=4pt, bottom=4pt,
  fonttitle=\bfseries\small, title=#1, coltitle=white}

% ---------------------------------------------------------------- code
\lstdefinelanguage{SV}{
  morekeywords={module,endmodule,input,output,wire,reg,signed,always,posedge,
    negedge,if,else,begin,end,assign,parameter,localparam,generate,endgenerate,
    for,genvar,case,endcase,function,endfunction,integer,initial},
  sensitive=true, morecomment=[l]{//}, morecomment=[s]{/*}{*/},
  morestring=[b]", morestring=[b]`}
\lstset{
  basicstyle=\ttfamily\small, backgroundcolor=\color{codebg},
  keywordstyle=\color{accent}\bfseries, commentstyle=\color{muted}\itshape,
  frame=single, rulecolor=\color{rule}, framesep=4pt, xleftmargin=4pt,
  xrightmargin=4pt, columns=fullflexible, keepspaces=true, breaklines=true,
  showstringspaces=false, upquote=true, aboveskip=0.6em, belowskip=0.6em}
\newcommand{\cmd}[1]{\texttt{#1}}
\newcommand{\file}[1]{\texttt{#1}}

%     \newcommand{\doctitle}{...}
%
% Team names live in ONE place (below). Edit them here and every document
% updates on the next build.
% ============================================================================

% ---------------------------------------------------------------- team
\newcommand{\projname}{FPGA int8 MLP Accelerator}
\newcommand{\repourl}{https://github.com/itsKannapi/fpga-mlp-accelerator}
\newcommand{\memberA}{Joey (\texttt{itsKannapi})}
\newcommand{\memberB}{Teammate~2 (TBD)}
\newcommand{\memberC}{Teammate~3 (TBD)}

% ---------------------------------------------------------------- packages
\usepackage[margin=1in]{geometry}
\usepackage[T1]{fontenc}
\usepackage{lmodern}
\usepackage{microtype}
\usepackage{amsmath}
\usepackage{booktabs}
\usepackage{tabularx}
\usepackage{longtable}
\usepackage{array}
\usepackage{enumitem}
\usepackage{xcolor}
\usepackage{listings}
\usepackage{tikz}
\usetikzlibrary{arrows.meta, positioning, fit, calc}
\usepackage[most]{tcolorbox}
\usepackage{fancyhdr}
\usepackage{titlesec}
\usepackage{needspace}
\usepackage[hidelinks]{hyperref}
\hypersetup{colorlinks=true, linkcolor=accent, urlcolor=accent, citecolor=accent}

% ---------------------------------------------------------------- colors
\definecolor{accent}{HTML}{1F4E79}
\definecolor{muted}{HTML}{5F6B7A}
\definecolor{rule}{HTML}{C9D1DA}
\definecolor{okc}{HTML}{2E7D32}
\definecolor{progc}{HTML}{B26A00}
\definecolor{todoc}{HTML}{8A8F98}
\definecolor{codebg}{HTML}{F5F7FA}

% ---------------------------------------------------------------- layout
\setcounter{tocdepth}{1}   % sections only, keeps the contents to a few lines
\setlength{\parindent}{0pt}
\setlength{\parskip}{0.55em}
\setlist{itemsep=0.15em, topsep=0.3em}
\renewcommand{\arraystretch}{1.18}

\titleformat{\section}{\Large\bfseries\color{accent}}{\thesection}{0.7em}{}
\titleformat{\subsection}{\large\bfseries}{\thesubsection}{0.6em}{}
\titleformat{\subsubsection}{\normalsize\bfseries}{\thesubsubsection}{0.5em}{}
\titlespacing*{\section}{0pt}{1.4em}{0.5em}

\pagestyle{fancy}
\fancyhf{}
\fancyhead[L]{\small\color{muted}\projname}
\fancyhead[R]{\small\color{muted}\doctitle}
\fancyfoot[C]{\small\color{muted}\thepage}
\renewcommand{\headrulewidth}{0.4pt}
\renewcommand{\headrule}{\color{rule}\hrule width\headwidth height\headrulewidth}

% ---------------------------------------------------------------- title block
\newcommand{\makedoctitle}[2]{%  #1 = subtitle, #2 = last-updated date
  \thispagestyle{plain}
  {\color{muted}\small\MakeUppercase{\projname}}\par\vspace{0.2em}
  {\Huge\bfseries\color{accent}\doctitle}\par\vspace{0.4em}
  {\large #1}\par\vspace{0.6em}
  {\small\color{muted}Team: \memberA, \memberB, \memberC\quad$\cdot$\quad
   Last updated #2\quad$\cdot$\quad\href{\repourl}{\texttt{\repourl}}}\par
  \vspace{0.4em}{\color{rule}\rule{\linewidth}{0.6pt}}\par\vspace{0.6em}
}

% ---------------------------------------------------------------- status tags
\newcommand{\tagbox}[2]{\tcbox[on line, colback=#1!12, colframe=#1!12,
  coltext=#1!70!black, boxsep=0pt, left=3pt, right=3pt, top=1.5pt, bottom=1.5pt,
  arc=2pt]{\footnotesize\bfseries #2}}
\newcommand{\done}{\tagbox{okc}{DONE}}
\newcommand{\inprog}{\tagbox{progc}{IN PROGRESS}}
\newcommand{\todo}{\tagbox{todoc}{TO DO}}
\newcommand{\stretchgoal}{\tagbox{accent}{STRETCH}}

% ---------------------------------------------------------------- callouts
\newtcolorbox{note}[1][Note]{colback=accent!5, colframe=accent!60,
  boxrule=0.6pt, arc=2pt, left=6pt, right=6pt, top=4pt, bottom=4pt,
  fonttitle=\bfseries\small, title=#1, coltitle=white}
\newtcolorbox{warn}[1][Watch out]{colback=progc!6, colframe=progc!70,
  boxrule=0.6pt, arc=2pt, left=6pt, right=6pt, top=4pt, bottom=4pt,
  fonttitle=\bfseries\small, title=#1, coltitle=white}

% ---------------------------------------------------------------- code
\lstdefinelanguage{SV}{
  morekeywords={module,endmodule,input,output,wire,reg,signed,always,posedge,
    negedge,if,else,begin,end,assign,parameter,localparam,generate,endgenerate,
    for,genvar,case,endcase,function,endfunction,integer,initial},
  sensitive=true, morecomment=[l]{//}, morecomment=[s]{/*}{*/},
  morestring=[b]", morestring=[b]`}
\lstset{
  basicstyle=\ttfamily\small, backgroundcolor=\color{codebg},
  keywordstyle=\color{accent}\bfseries, commentstyle=\color{muted}\itshape,
  frame=single, rulecolor=\color{rule}, framesep=4pt, xleftmargin=4pt,
  xrightmargin=4pt, columns=fullflexible, keepspaces=true, breaklines=true,
  showstringspaces=false, upquote=true, aboveskip=0.6em, belowskip=0.6em}
\newcommand{\cmd}[1]{\texttt{#1}}
\newcommand{\file}[1]{\texttt{#1}}

\newcommand{\doctitle}{Developer Setup Guide}

\begin{document}
\makedoctitle{From a fresh machine to passing testbenches, and how we work together}{2026-10-16}

\begin{note}[Short version]
Clone the repo $\rightarrow$ install Vivado 2026.1 (Artix-7 only) and its free
Basic license $\rightarrow$ in Vivado's Tcl console,
\cmd{source <repo>/scripts/create\_project.tcl} $\rightarrow$ Run Simulation
$\rightarrow$ look for \texttt{ALL TESTS PASSED}. No Vivado (e.g.\ on a Mac)? Run
\cmd{./scripts/run\_sims.sh} with Icarus Verilog instead.
\end{note}

\tableofcontents

% ============================================================================
\section{What you need}
\begin{center}
\begin{tabularx}{\linewidth}{@{}lX@{}}
\toprule
Item & Notes \\
\midrule
Windows or Linux PC & \textbf{Vivado does not run on macOS.} Mac users can still run
  simulations with Icarus (Section~\ref{sec:icarus}) and do Python and docs work. \\
16\,GB RAM, $\sim$60\,GB disk & For Vivado with a single device family installed. \\
Git & Windows: \cmd{winget install Git.Git} \\
Python 3.12+ & Only needed to retrain or re-export weights. 3.14 works. \\
Basys~3 + micro-USB cable & Only for hardware steps. All simulation works without the board. \\
\bottomrule
\end{tabularx}
\end{center}

% ============================================================================
\section{Get the code}
\begin{lstlisting}[language=bash]
git clone https://github.com/itsKannapi/fpga-mlp-accelerator.git
cd fpga-mlp-accelerator
git config --global user.name  "Your Name"
git config --global user.email "ID+username@users.noreply.github.com"
\end{lstlisting}
Use the \texttt{noreply} address from GitHub $\rightarrow$ Settings $\rightarrow$
Emails (``Keep my email addresses private''). Every commit's email is public in a
public repo. Commits still count on your profile.

Clone somewhere with a short path with no spaces (e.g.\ \file{C:\textbackslash dev\textbackslash}),
not a OneDrive-synced folder: OneDrive tries to sync Vivado's gigabytes of build files.

% ============================================================================
\section{Install Vivado 2026.1}
\begin{enumerate}
  \item Download the AMD unified installer (free AMD account required).
  \item Choose \textbf{Vivado}, Standard/Basic edition. In the device list,
        \textbf{deselect everything except 7 Series $\rightarrow$ Artix-7}. Skip Vitis.
        This keeps the install near 60\,GB instead of 100\,GB or more.
  \item If you later search for part \texttt{xc7a35tcpg236-1} and get 0 matches, Artix-7
        wasn't installed: \emph{Help $\rightarrow$ Add Design Tools or Devices}
        (or re-run \file{xsetup.exe}) and add it.
\end{enumerate}

\subsection{License (required even for the free tier)}
From 2026.1, Vivado will not start without a license file, including the free
tier. The error is \emph{``Vivado Design Suite cannot be launched because a
valid license was not found.''}
\begin{enumerate}
  \item Open the Vivado License Manager $\rightarrow$ \emph{View System Information}.
        Write down a Host ID (MAC address or disk serial).
  \item On AMD's licensing site, select \textbf{``Vivado Basic Tier License, Node
        Locked License''} (Certificate, No Charge).
  \item Lock it to the Host ID from step 1, \textbf{character for character}.
        A mismatched Host ID is the most common cause of the error above.
  \item License Manager $\rightarrow$ \emph{Load License} $\rightarrow$ select the
        \file{.lic} file. Check \emph{View License Status} shows
        \textbf{HostID matches = YES}.
\end{enumerate}
\begin{warn}[Do not pick these]
\textbf{``Start Now! 60 Day Trial''} and \textbf{``Vivado Enterprise Tier, 60 Day
Evaluation''} expire mid-project. \textbf{``HL WebPACK 2015 and Earlier''} only
works for old tool versions. The Basic license renews yearly.
\end{warn}

% ============================================================================
\section{Create the Vivado project}\label{sec:proj}
The project folder \file{vivado/} is not in git. It is generated, and it would
contain absolute paths from whoever made it. Build your own:
\begin{enumerate}
  \item Start Vivado. At the start screen, open the \textbf{Tcl Console} tab at the bottom.
        (If a project is open, \emph{File $\rightarrow$ Close Project} first.)
  \item Run, with forward slashes:
\begin{lstlisting}[language=tcl]
source C:/dev/fpga-mlp-accelerator/scripts/create_project.tcl
\end{lstlisting}
  \item It prints \texttt{Project ready: ...mlp\_accel.xpr}. The project is open.
\end{enumerate}
The script writes \file{vivado/data\_dir.vh} (where \file{rtl/} lives on
\emph{your} machine), adds every \file{src/*.v} and \file{sim/*.v}
\textbf{without copying them}, sets the include paths, and sets the simulation to
run to completion. Re-run it any time the project gets into a weird state, or
after a teammate adds a new source file.

% ============================================================================
\section{Run the testbenches}

\subsection{In Vivado}
\emph{Flow Navigator $\rightarrow$ Run Simulation $\rightarrow$ Run Behavioral
Simulation}. Read the result in the Tcl Console (scroll up if needed). To run a
different testbench, type
\cmd{set\_property top tb\_mac [get\_filesets sim\_1]} and run again.

\begin{center}
\begin{tabular}{@{}lll@{}}
\toprule
Testbench & Expected last lines & Time in xsim \\
\midrule
\file{tb\_mac}    & \texttt{ALL TESTS PASSED (2266 checks)} & $\sim$20\,s \\
\file{tb\_layer1} & \texttt{ALL TESTS PASSED (4259 checks)} & under a minute \\
\bottomrule
\end{tabular}
\end{center}
\file{tb\_layer1} also prints \emph{exercised ReLU 41 times, clamp 2 times}. If
you see \texttt{could not open golden.txt}, the project wasn't made by the script:
re-run Section~\ref{sec:proj}.

\subsection{Without Vivado: Icarus Verilog}\label{sec:icarus}
Fast, free, runs on macOS, Linux, and Git Bash on Windows. Use it as a quick
check before every push.
\begin{lstlisting}[language=bash]
brew install icarus-verilog        # macOS
sudo apt install iverilog          # Ubuntu
./scripts/run_sims.sh
\end{lstlisting}
It runs every \file{sim/tb\_*.v}, prints one line per testbench, and exits
non-zero if anything fails.

% ============================================================================
\section{Python (only to retrain or re-export)}
\begin{lstlisting}[language=bash]
python -m venv .venv              # Windows: py -m venv .venv
.venv\Scripts\activate            # macOS/Linux: source .venv/bin/activate
pip install torch torchvision numpy
python quantize_export.py
\end{lstlisting}
\begin{warn}[Weights are frozen]
Don't re-run \file{train\_mlp.py}: retraining gives different weights and
breaks every golden comparison. If \cmd{python} opens the Microsoft Store on
Windows, use \cmd{py}, or turn off the \texttt{python.exe} alias under
\emph{Settings $\rightarrow$ Apps $\rightarrow$ Advanced app settings $\rightarrow$
App execution aliases}. Re-exporting currently needs a file only Joey has
(open issue \#1 in \file{progress.tex}).
\end{warn}

% ============================================================================
\section{Programming the board}
\begin{enumerate}
  \item Plug into the \textbf{PROG} USB port and turn on the board's power switch.
  \item \emph{Open Hardware Manager $\rightarrow$ Open Target $\rightarrow$ Auto Connect}.
  \item Click \textbf{Program Device} and select the \file{.bit} file. Opening the
        target is not enough: the board shows ``Programmed'' even while it is still
        running the factory demo from flash.
\end{enumerate}
Programming over JTAG is lost at power-off; reprogram each session. If Auto Connect
finds nothing: check the cable carries data (not charge-only), try a port on the
motherboard, or install Digilent Adept for its USB drivers.

% ============================================================================
\section{How we work together}
\begin{enumerate}
  \item \textbf{Branch per task}: \cmd{git switch -c layer2-argmax}. Don't commit
        straight to \texttt{main}.
  \item \textbf{Every new module comes with a self-checking testbench} that prints
        \texttt{ALL TESTS PASSED}, and a few planted bugs to prove it can fail
        (\file{design.tex} \S4).
  \item \textbf{Before opening a pull request}: \cmd{./scripts/run\_sims.sh} passes,
        \file{progress.tex} has a dated entry, and the PDFs are rebuilt.
  \item \textbf{Pull request}: push the branch, open a PR on GitHub, and have one
        teammate review it before merging into \texttt{main}.
  \item \textbf{Never commit} \file{vivado/}, \file{.venv/}, \file{data/}, or
        \file{*.npz}. \file{.gitignore} already blocks them; check
        \cmd{git status} before \cmd{git add}.
  \item \textbf{Commit messages} say what changed and how it was verified, e.g.\
        \emph{``Add layer 2 + argmax; bit-exact vs golden on 64 images''}.
\end{enumerate}

\subsection{Building these documents}
\begin{lstlisting}[language=bash]
cd docs
latexmk -pdf design.tex progress.tex setup.tex
\end{lstlisting}
Needs a LaTeX distribution: MiKTeX on Windows, MacTeX on macOS, \texttt{texlive}
on Linux. Or upload \file{docs/} to Overleaf. The PDFs are committed so they can
be read on GitHub; rebuild and commit them whenever a \file{.tex} changes.

\end{document}
